\begin{abstract}
现有时间序列异常检测基准回答“一个分数能否预测异常”，却没有回答“它在简单统计信号之外增加了什么”。本文提出条件检测效用（CDU）：加入检测器分数后，统计参考的留出对数损失降低量。我们采用留一来源评价，在覆盖 23 个来源组的 350 条单变量序列上比较九个检测器的固定异常分数。在相同 spline 探针、损失和来源权重下，MOMENT-FT 与 MOMENT-ZS 的检测器独立效用分别排第 4、5 名，CDU 却排第 7、9 名；TranAD 从第 6 升至第 4。POLY 与 MOMENT-ZS 的基准 VUS-PR 几乎相同（0.389 对 0.383），CDU 却相差 0.00593 比特。MOMENT-ZS 的分数可被统计参考高度重构（$R^2=0.904$），CDU 却排第 9；九模型的可重构性与 CDU 排名对应很弱（Spearman $\rho=-0.167$）。三类探针保持相同前两名次序和前四名集合；五个抽样种子得到相同 spline 排名，参考变体也保留领先组。独立表现与参考之外的贡献是实证上不可互换的评价轴；CDU 使后者可在异构 TSAD 分数上被测量。
\end{abstract}
